% !TEX root = ../notes.tex

\documentclass[../notes.tex]{subfiles}

\begin{document}

\section{April 1}
Here we go.

\subsection{Algebraic Geometry}
Our current goal is to find a Galois extension of $\mathbb E_\infty$-rings $A\to B$ where $A$ is something interesting, and $B$ is computable (e.g., an $\mathbb E_\infty$-ring).
\begin{example}
	Last class we found that $L_{K(1)}\mathbb S=(\mathrm{KU}_p^\land)^{h\ZZ_p^\times}$. Let's quickly recall what $\mathrm{KU}^\land_p$ is.
	\begin{itemize}
		\item We could realize $\mathrm{KU}^\land_p$ as an $\mathbb F_p^\times$-Galois extension of $(\mathrm{BP}/(v_2,v_3,\ldots))^\land_p[1/v_1]$.
		\item Alternatively, $\mathrm{KU}_p^\land=\mathrm{ku}^\land_p[1/\beta]$, where $\mathrm{ku}$ is the group completion of the category of vector spaces.
		\item Lastly, we could take $\mathrm{KU}^\land_p=(\Sigma_+^\infty\CP^\infty)^\land_p[1/\beta]$.
	\end{itemize}
\end{example}
In particular, we will attempt to compute with $L_{K(n)}\mathbb S$ to imitate the construction in the third point and then use something like the first point. Namely, we will take some complicated spaces and invert some elements to hopefully get even $\mathbb E_\infty$-rings.
\begin{example}[cyclotomic]
	There are Galois extensions of $L_{K(n)}\mathbb S$ which look like
	\[L_{K(n)}\Sigma_+^\infty K(A,n)[1/z].\]
\end{example}
Our goal will be to construct a functor $E(R_0,\Gamma_0)$ which produces a universal deformation, which is an $\mathbb E_\infty$-ring equipped with a formal group deforming $\Gamma_0$. Accordingly, we should start by explaining what it means to have a formal group over an $\mathbb E_\infty$-ring, which requires a little spectral algebraic geometry.
\begin{definition}[connective]
	An $\mathbb E_\infty$-ring $R$ is \textit{connective} if and only if $\pi_iR=0$ for negative $i$.
\end{definition}
\begin{definition}[spectrum]
	Given a connective $\mathbb E_\infty$-algebra $A$ over a connective $\mathbb E_\infty$-ring $R$. Then we define its \textit{spectrum} $\Spec A$ as the functor taking in connective $\mathbb E_\infty$-algebras $B$ over $R$ to $\Spec A(B)\coloneqq\op{Hom}_{\mathbb E_\infty}(A,B)$.
\end{definition}
\begin{remark}
	It turns out to be better to work with test objects which are even $\mathbb E_\infty$-algebras. (Roughly speaking, the functor $\Spec A$ on connective algebras can be recovered from the functor which only takes even algebras as input.) This has to do with motivic homotopy theory.
\end{remark}
\begin{definition}[reduction]
	If $A$ is connective, then we define $A^{\mathrm{red}}\coloneqq(\pi_0A)^{\mathrm{red}}$.
\end{definition}
Notably, the reduction forgets about all the homotopy information!
\begin{example}
	One can imagine taking the intersection of a point on a line with itself as a tensor product $\CC\otimes_{\CC[x]}\CC$. This has homotopy groups concentrated in the two dimensions $0$ and $1$, each of one complex dimension, which recovers a notion of multiplicity. However, the actual geometric object is the single point living in $\pi_0$.
\end{example}
Here are some examples of spectra.
\begin{example}
	We define $\mathbb G_m$ over $R$ to be $\Spec(R\otimes\Sigma^\infty_+\ZZ)$. To convince us that this is a reasonable definition, we note that $\pi_*(R\otimes\Sigma_+^\infty\ZZ)=\pi_*(R)[x,1/x]$, where $x$ lives in degree zero. Its functor of points takes ca connective algebra $A$ to
	\begin{align*}
		\mathbb G_m(A) &= \op{Hom}_{\op{Alg}_R(\mathbb E_\infty)}(R\otimes\Sigma_+^\infty\ZZ,A) \\
		&= \op{Hom}_{\mathrm{Alg}_\ZZ(\mathbb E_\infty)}(\Sigma_+^\infty\ZZ,A) \\
		&= \op{Hom}_{\mathrm{Spectra}}(\ZZ,\op{GL}_1A),
	\end{align*}
	where we are using the fact that $\Sigma_+^\infty$ mapping connective grouplike $\mathbb E_\infty$-spaces to $\mathbb E_+^\infty$-rings admits a right adjoint called $\op{GL}_1$. (This is like taking the group of units of a ring.) Notably, $\op{GL}_1$ is $\Omega^\infty$ of some ring $\op{gl}_1$, so we can factor out a copy of $\Omega^\infty$. The fact that this outputs to spectra instead of spaces is analogous to the fact that $\mathbb G_m$ classically outputs to groups instead of sets. %In fact, it follows that $\mathbb G_m$ outputs to $\mathrm{Mod}(\ZZ)$.
\end{example}
% \begin{remark}
% 	In fact, $\mathbb G_m$ outputs $\ZZ$-modules, which one could see by our last calculation.
% \end{remark}
\begin{remark}
	Note that $\mathbb G_m$ fits into the adjunction
	\[\op{Hom}_{\mathrm{Mod}(\ZZ)}(M,\mathbb G_m(A))=\op{Hom}_{\mathbb E_\infty}(\Sigma_+^\infty\Omega^\infty M,A).\]
	One could take this as a definition of $\mathbb G_m$, which shows that $\mathbb G_m$ outputs to $\mathrm{Mod}(\ZZ)$.
\end{remark}
\begin{example}
	We define $\widehat{\mathbb G}_m$ to be $\op{Spf}R[x,1/x]^\land_{(1-x)}$. (Note that this completion can be seen to be an $\mathbb E_\infty$-ring because it is Bousfield localization for the module $\pi_0A/(1-x)$.) Alternatively, we could define this as the colimit of $\Spec R[x,1/x]/(1-x)^n$. Formally, the $A$-points are given by the subspace of
	\[\op{Hom}_{\mathrm{Alg}_R(\mathbb E_\infty)}(R[x,1/x],A),\]
	where the image of $1-x$ is nilpotent. Here, $\pi_*R[x,1/x]^\land_{1-x}=\pi_*R[[y]]$ (one can show that $\pi_*$ passes through completions, with some elbow grease), where $y=1-x$. It again turns out that $\widehat{\mathbb G}_m$ outputs to $\mathrm{Mod}(\ZZ)$, but our starting presentation did not know this.
\end{example}
\begin{remark}
	It turns out that $R[y]/y$ is always an $\mathbb E_\infty$-ring. In general, it also turns out that $R/y^N$ is an $\mathbb E_k$-ring, where $k$ grows with $N$.
\end{remark}

\subsection{Formal Groups}
The point of the previous example is to give us some examples of formal groups.
\begin{definition}[formal group]
	Fix a connective $\mathbb E_\infty$-ring $R$. Then a \textit{formal group} over $R$ is a functor $G$ from connective $R$-algebras to connective $\ZZ$-modules such that $\Omega^\infty\circ G$ is a formal hyperplane.
\end{definition}
Wait, what is a formal hyperplane?
\begin{example}
	We see that $\op{Spf}R[[y]]$ is a formal hyperplane, as is $\op{Spf}R[[y_1,\ldots,y_n]]$. In general, formal hyperplanes all locally look like these.
\end{example}
Ok, so what is a formal hyperplane? Well, such a thing should be flat.
\begin{definition}[flat]
	Fix a connective $\mathbb E_\infty$-ring $R$. Then an $R$-module $M$ is \textit{flat} if and only if $\pi_0M$ is flat over $\pi_0R$, and $\pi_iM\cong\pi_0M\otimes_R\pi_iR$ for all $i$.
\end{definition}
\begin{remark}
	This is a very strong notion of flatness: indeed, it permits little interesting homotopy theory. There are more general definitions of flatness, but we will not need them.
\end{remark}
\begin{example}
	A module $M$ over a skew field $R$ is flat if and only if $M$ is a direct sum of copies of $R$. Notably, there are examples of modules over skew fields which are not flat!
\end{example}
We are now really ready to define a formal hyperplane.
\begin{definition}[formal hyperplane]
	Fix a connective $\mathbb E_\infty$-ring $R$. A \textit{formal hyperplane} is $\op{Spf}$ of the $R$-dual of a co-free $\mathbb E_\infty$-coalgebra of a flat $R$-module.
\end{definition}
% \begin{remark}
% 	It is equivalent to ask for the dual of a projective $R$-module, meaning that it is a retract of a free module.
% \end{remark}
\begin{example}
	Note that $R[[y]]=\op{Hom}_R(R[y],R)$, so $\op{Spf}R[[y]]$ is a formal hyperplane.
\end{example}
Here is another way to view this definition: if $M$ is a discrete module over a discrete ring $R$, then we may define $\op{Sym}M\coloneqq\bigoplus_i\op{Sym}^iM$ and its completion $\op{Sym}^\land M\coloneqq\prod_i\op{Sym}^iM$. Note that $\op{Sym}^\land M$ is the linear dual of the co-free algebra on $M$. Moving back to homotopy theory, here is the corresponding definition.
\begin{definition}[smooth]
	Fix an $\mathbb E_\infty$-ring $R$. Then a \textit{smooth coalgebra} is an $\mathbb E_\infty$-algebra $X$ is one which is flat over $R$ and such that $\pi_0X$ is a co-free coalgebra on a projective $R$-module. Here, projective means that the module is a retract of a free module.
\end{definition}
Then taking $\op{Spf}\op{Hom}_R(X,R)$ of a smooth coalgebra $X$ again recovers the notion of a formal hyperplane.

We are now allowed to give some examples of formal groups.
\begin{example}
	Over $\ZZ$, we see that $\widehat{\mathbb G}_m$ is $\op{Spf}\ZZ[[y]]$, and its group structure is given by the morphism $\ZZ[[y]]\to\ZZ[[y_1,yy_2]]$ given by $y\mapsto(y_1+y_2+y_1y_2)$.
\end{example}
\begin{example}
	Over $\ZZ$, we may define $\widehat{\mathbb G}_a$ as $\op{Spf}\ZZ[[y]]$ with the formal group law given by the morphism $\ZZ[[y]]\to\ZZ[[y_1,y_2]]$ given by $y\mapsto y_1+y_2$.
\end{example}
We will need to be careful about our deformations, as the following example shows.
\begin{exe}
	There is no formal group over $\tau_{\le1}\mathbb S$ which base-changes to $\widehat{\mathbb G}_a$ over $\ZZ=\tau_{\le0}\mathbb S$
\end{exe}
\begin{proof}
	Suppose that there is such a deformation $G=\op{Spf}A$. Note that $A\otimes_{\tau_{\le1}\mathbb S}\ZZ=\ZZ[[y]]$, so $A$ looks something like $\tau_{\le1}\mathbb S[[x]]$, which we do not expect to be a formal group because the element $\eta$ has strange properties with distributivity. To be formal, note that $\pi_1A=\pi_0A/2$, induced by the map $\eta\colon\pi_0A\to\pi_1A$ (which has $2\eta=0$). The isomorphism $A\otimes_{\tau_{\le1}\mathbb S}\ZZ=\ZZ[[y]]$ also (by K\"unneth) produces isomorphisms $\alpha$ and $\beta$ living in the commutative diagram
	% https://q.uiver.app/#q=WzAsNixbMCwwLCJcXFpaW1t5XV0iXSxbMCwxLCJcXFpaW1t5XzEseV8yXV0iXSxbMSwwLCJcXHBpXzBBIl0sWzEsMSwiXFxwaV8wKEFcXG90aW1lcyBBKSJdLFsyLDAsIlxccGlfMUEiXSxbMiwxLCJcXHBpXzEoQVxcb3RpbWVzIEEpIl0sWzIsNCwiXFxldGEiXSxbMyw1LCJcXGV0YSJdLFs0LDVdLFsxLDNdLFswLDJdLFsyLDNdLFswLDFdXQ==&macro_url=https%3A%2F%2Fraw.githubusercontent.com%2FdFoiler%2Fnotes%2Fmaster%2Fnir.tex
	\[\begin{tikzcd}[cramped]
		{\ZZ[[y]]} & {\pi_0A} & {\pi_1A} \\
		{\ZZ[[y_1,y_2]]} & {\pi_0(A\otimes A)} & {\pi_1(A\otimes A)}
		\arrow[from=1-1, to=1-2]
		\arrow[from=1-1, to=2-1]
		\arrow["\eta", from=1-2, to=1-3]
		\arrow[from=1-2, to=2-2]
		\arrow[from=1-3, to=2-3]
		\arrow[from=2-1, to=2-2]
		\arrow["\eta", from=2-2, to=2-3]
	\end{tikzcd}\]
	where the left two vertical arrows are the formal group law.
	
	To continue, we need to use something about $\mathbb E_\infty$-space structure. Now, if $X$ is any $\mathbb E_\infty$-space with $x\in\pi_0X$, then note that $x$ factors through $\op{Free}_{\mathbb E_\infty}\mathrm{pt}=*\sqcup\mathrm B\Sigma_2\sqcup\mathrm B\Sigma_3\sqcup\cdots$, so we receive a class in $\pi_1(X)$ which is the image from $\pi_1(\mathrm B\Sigma_2)$; if $X$ comes from a spectrum (by $\Omega^\infty$), then one can calculate that this map $\pi_0X\to\pi_1X$ is given by $\eta$.

	Now, $A$ is an $\mathbb E_\infty$-ring, so $\Omega^\infty A$ has two structures as an $\mathbb E_\infty$-space (namely, an additive and a multiplicative one). Applying this to the above construction to these two operations, we receive two maps $\eta_a,\eta_m\colon\pi_0A\to\pi_1A$; note that $\eta_a$ is the map $\eta\colon\pi_0A\to\pi_1A$ given earlier. By fiddling with the distributive law, we see that
	\[\eta_m(x_1+x_2)=\eta_mx_1+\eta_mx_2+\eta_ax_1x_2.\]
	(This is a fancy version of $(x_1+x_2)^2=x_1^2+x_2^2+2x_1x_2$.) On the other hand, there is $f\in\ZZ[[t]]$ (well-defined up to$\pmod2$) for which $\eta_m\alpha(t)=\eta f(t)$ because $\eta\colon\pi_0A\to\pi_1A$ is an isomorphism, so tracking around the diagram from earlier show that
	\[\eta_m\beta(t_1+t_2)=\eta\beta f(t_0+t_1).\]
	One can expand out the left-hand side to see that the coefficient of $t_0t_1$ is $\eta$, but it should be even (because $f$ is just some power series).
\end{proof}

\end{document}